\chapter{Arquitectura y estructura del proyecto}\label{cap:arquitectura}

\section{Visión general}

El proyecto tiene dos partes que se comunican por archivos:

\begin{enumerate}
  \item \textbf{Datos} (\texttt{analisis/}): scripts SQL de DuckDB que cargan los CSV de
        FIRMS, calculan las constantes, armonizan cada foco y exportan los productos de
        la sección~\ref{sec:contrato}.
  \item \textbf{Aplicación} (\texttt{frontend/}): página web que lee esos productos como
        archivos estáticos y hace todos los cálculos de área y periodo en el navegador.
\end{enumerate}

Flujo de los datos:

\begin{center}
\small
CSV de FIRMS $\xrightarrow{\text{01\_cargar}}$
\texttt{focos.duckdb} $\xrightarrow{\text{02, 03}}$
constantes $k$ $\xrightarrow{\text{04\_armonizar}}$
focos con peso $\xrightarrow{\text{05\_exportar}}$
\texttt{frontend/public/datos/} $\to$ navegador
\end{center}

\section{Estructura del repositorio}

\begin{verbatim}
Hots-Spots-Nasa-Challenge/
|-- README.md              Reto, problema y criterios de éxito
|-- mise.toml              Herramientas (versiones) y tareas del proyecto
|-- compose.yaml           Entorno Docker: frontend + backend
|-- analisis/              Parte de datos (DuckDB)
|   |-- 01_cargar.sql          CSV -> tabla focos
|   |-- 02_correlacion.sql     Correlación MODIS-VIIRS
|   |-- 03_k_mensual.sql       Constantes por mes
|   |-- 04_armonizar.sql       Peso por foco y serie mensual
|   |-- 05_exportar_mapa.sql   Rejilla 0,1° x mes para el mapa
|   |-- *.md                   Resultados y validación
|   `-- focos.duckdb           Base local (no se versiona)
|-- frontend/              Aplicación web
|   |-- src/
|   |   |-- App.tsx            Pantalla principal y controles
|   |   |-- Mapa.tsx           Mapa (MapLibre GL)
|   |   `-- datos.ts           Carga y agregación de los datos
|   `-- public/datos/      Productos exportados (JSON)
|-- backend/               API FastAPI (esqueleto, solo /salud)
`-- Documentacion/         Este documento y los planes
    |-- main.tex               Documento de requerimientos
    |-- NN-capitulo/           Un capítulo por subcarpeta
    |-- avance.tex             Estado del proyecto
    `-- frontend/plan.tex      Propuesta inicial del frontend
\end{verbatim}

Los CSV crudos viven fuera del repo, en \texttt{../data de prueba/}.

\section{Tecnologías}

\begin{tabularx}{\linewidth}{@{}lX@{}}
\toprule
Pieza & Elección y motivo \\
\midrule
Datos & DuckDB: SQL sobre millones de filas en un portátil, sin servidor. \\
Frontend & React 19 + TypeScript + Vite + Tailwind 4. \\
Mapa & MapLibre GL JS con el mapa base \emph{Positron} de OpenFreeMap (claro, sin
       clave de API). \\
Consultas en el navegador & Hoy, JSON agregado en JavaScript. Previsto: Parquet +
       DuckDB-WASM para áreas y periodos arbitrarios. \\
Gráficos & SVG propio para el calendario (es una cuadrícula). \\
Backend & FastAPI + DuckDB, solo si hace falta (ver decisiones abiertas). \\
Entorno & \texttt{mise} (herramientas y tareas), Docker Compose / Podman,
          \texttt{uv} para Python, \texttt{pnpm} para Node, Tectonic para \LaTeX. \\
\bottomrule
\end{tabularx}

\section{Pantalla principal}

Una sola pantalla, de arriba abajo:

\begin{enumerate}
  \item Barra superior mínima: nombre, región, ayuda.
  \item Selector de periodo (año y mes).
  \item Mapa con la herramienta de selección de área.
  \item Frase de respuesta (anomalía y ranking del periodo).
  \item Calendario años~$\times$~meses con su leyenda.
  \item Rankings y descarga, discretos, al final.
\end{enumerate}

\abierta{si el mapa ocupa la mayor parte de la pantalla con el calendario como panel
lateral, o si el calendario es el protagonista y el mapa un selector pequeño.}

\section{Tareas de \texttt{mise}}

\begin{tabularx}{\linewidth}{@{}lX@{}}
\toprule
Tarea & Qué hace \\
\midrule
\texttt{datos:cargar} & Carga los CSV en \texttt{analisis/focos.duckdb}. \\
\texttt{datos:correlacion} & Correlación y constantes de conversión. \\
\texttt{datos:armonizar} & Serie armonizada 2000--2025. \\
\texttt{datos:mapa} & Exporta la rejilla para el mapa. \\
\texttt{front:dev} / \texttt{front:build} & Servidor de desarrollo / build de producción. \\
\texttt{back:dev} & API en modo desarrollo. \\
\texttt{docker:up} / \texttt{docker:down} & Entorno completo en contenedores. \\
\texttt{docs:pdf} & Compila los \texttt{.tex} de \texttt{Documentacion/}. \\
\bottomrule
\end{tabularx}
